% Part: propositional-logic
% Chapter: propositional-logic
% Section: completeness
%READERNOTE{SOL6-A338: द्विदिशा दर्शवणारा वाक्प्रयोग आणि विगमनाची संज्ञा सुसंगत केली; संक्षिप्त सिद्धता आणि सराव जपले.}

\documentclass[../../../include/open-logic-section]{subfiles}

\begin{document}

\olfileid{pl}{prp}{com}

\olsection{विधानात्मक तर्कशास्त्राची संपूर्णता}

\begin{defn}
\ollabel{def:MaxCon}
!!{formula}s चा संच~$\Gamma$ \emph{महत्तम सुसंगत} असतो, जर तो
सुसंगत असेल आणि $\Gamma \subseteq \Delta$ असा कोणताही सुसंगत
संच~$\Delta$ घेतल्यास $\Gamma = \Delta$ असेल.
\end{defn}

\begin{lem}[सत्यता पूर्वप्रमेय]
  \ollabel{lem:truth}
$\Gamma$ महत्तम सुसंगत असू द्या; मग:
\begin{enumerate}
\item \ollabel{lem:truth-1} $!A \in \Gamma$ तेव्हा आणि केवळ तेव्हा $\lnot !A
  \notin \Gamma$;
\item \ollabel{lem:truth-2} $\Gamma \Proves !A$ तेव्हा आणि केवळ तेव्हा $!A \in \Gamma$;
\item \ollabel{lem:truth-3} $!A \lif !B \in \Gamma$ तेव्हा आणि केवळ तेव्हा
  $!A \notin \Gamma$ किंवा $!B \in \Gamma$.
\end{enumerate}
\end{lem}

\begin{proof}
\begin{enumerate}
\item $!A \in \Gamma$ असे माना. $\lnot !A \in \Gamma$ सुद्धा असेल,
  तर $\Gamma$ विसंगत ठरतो. $!A$ आणि $\lnot!A$ यांपैकी एकही
  $\Gamma$ मध्ये नसेल, तर \olref[axd]{prop:phi} नुसार
  $\Gamma\cup\{!A\}$ आणि $\Gamma \cup \{\lnot !A\}$ यांपैकी
  एक संच सुसंगत असतो; म्हणून $\Gamma$ महत्तम नसतो. उलटी दिशाही
  याचप्रमाणे सिद्ध होते.
\item $!A \in \Gamma$ असेल, तर अर्थातच $\Gamma \Proves !A$.
  उलट, $\Gamma \Proves !A$ असे माना. $!A \notin \Gamma$ असेल,
  तर \olref{lem:truth-1} नुसार $\lnot !A \in \Gamma$. त्यामुळे
  $\Gamma \Proves !A$ आणि $\Gamma \Proves \lnot !A$, हा व्याघात.
\item सरावासाठी.
\end{enumerate}
\end{proof}

\olref{lem:truth}\olref{lem:truth-2} ची डावीकडून उजवीकडील
दिशा दाखवते की $\Gamma$ \emph{निगमनतः बंद} आहे: म्हणजे
कोणत्याही~$!A$ साठी $\Gamma \Proves !A$ असेल, तर $!A
\in \Gamma$.

\begin{prob}
  \olref{lem:truth} ची सिद्धता पूर्ण करा.
\end{prob}

\begin{thm}[संपूर्णता]
\ollabel{thm:completeness}
$\Gamma$ सुसंगत असेल, तर $\Gamma$ पूर्ततायोग्य असतो.
\end{thm}

\begin{proof}
$!A_0$, $!A_1$,~\dots ही भाषेतील सर्व !!{formula}s ची
पूर्ण यादी असू द्या. पुढीलप्रमाणे !!{formula}s च्या संचांची
$\Gamma_0$, $\Gamma_1$,~\dots ही वाढती क्रमिका
आवर्ती रीतीने परिभाषित करू:
\begin{align*}
\Gamma_0 & = \Gamma\\
\Gamma_{n+1} &  =
\begin{cases}
  \Gamma_n \cup \{!A_n\} & \text{जर $\Gamma_n \cup \{!A_n\}$ सुसंगत असेल;}\\
  \Gamma_n \cup \{\lnot !A_n\} & \text{अन्यथा}.
\end{cases}
\end{align*}
मग पुढील व्याख्या करू:
\[
\Gamma^* = \bigcup_{0\le n}\Gamma_n.
\]
आता क्रमाने पुढील तथ्ये सिद्ध करायची आहेत:
\begin{enumerate}
\item प्रत्येक $n$ साठी $\Gamma_n$ हा संच सुसंगत आहे
  ($n$ वर विगमनाने, \olref[axd]{prop:phi} वापरून);
\item $\Gamma^*$ सुसंगत आहे;
\item $\Gamma^*$ महत्तम आहे.
\end{enumerate}
मग $p_i \in \Gamma^*$ असेल तेव्हाच $v(p_i) = \True$
असे ठरवून सत्य-मूल्यांकन~$v$ परिभाषित करू. $!A$ वर विगमन
केल्यास अधिक गुंतागुंतीच्या !!{sentence}s साठीही
$\Gamma^*$ मधील सदस्यत्व आणि $v$ नुसार सत्यता जुळतात,
हे दिसते: $\overline{v}(!A) = \True$ तेव्हा आणि केवळ तेव्हा
$!A \in \Gamma^*$. विशेषतः, $v$ हे $\Gamma^*$ ची पूर्तता
करते. $\Gamma \subseteq \Gamma^*$ असल्याने $\Gamma$
सुद्धा पूर्ततायोग्य आहे, हेच सिद्ध करायचे होते.
\end{proof}

\begin{cor}
$\Gamma \Entails !A$ असेल, तर $\Gamma \Proves !A$.
\end{cor}

\begin{proof}
$\Gamma\Proves/ !A$ असेल, तर \olref[axd]{prop:prov-incons} नुसार
$\Gamma \cup \{\lnot !A\}$ सुसंगत आहे. म्हणून संपूर्णता
प्रमेयानुसार $\Gamma \cup \{\lnot !A\}$ पूर्ततायोग्य आहे;
त्यामुळे $\Gamma \Entails/ !A$.
\end{proof}

\begin{prop}[संहतता प्रमेय]
\ollabel{prop:compactness}
$\Gamma$ पूर्ततायोग्य असेल तेव्हा आणि केवळ तेव्हा $\Gamma$ चा प्रत्येक
\emph{सान्त} उपसंच~$\Gamma_0$ पूर्ततायोग्य असेल.
\end{prop}

\begin{proof}
$\Gamma$ अपूर्ततायोग्य असेल तेव्हा आणि केवळ तेव्हा तो विसंगत असेल;
तेव्हा आणि केवळ तेव्हा $\Gamma$ चा एखादा सान्त उपसंच~$\Gamma_0$
विसंगत असेल; आणि तेव्हा आणि केवळ तेव्हा $\Gamma$ चा एखादा सान्त
उपसंच~$\Gamma_0$ अपूर्ततायोग्य असेल.
\end{proof}

\begin{cor}
$\Gamma \Entails !A$ असेल तेव्हा आणि केवळ तेव्हा $\Gamma$ च्या
एखाद्या सान्त उपसंच~$\Gamma_0$ साठी $\Gamma_0 \Entails !A$ असेल.
\end{cor}

\end{document}
